\chapter*{摘要}

撰写高质量的开题报告是博士研究生培养过程中的重要环节。一份规范、美观的开题报告不仅能够清晰地展示研究思路和计划，还能体现研究者的学术素养。随着计算机排版技术的发展，\LaTeX 因其强大的公式排版能力和优秀的文档管理功能，已成为学术界广泛使用的文档撰写工具。然而，上海创智学院的博士生开题报告目前只有基于Word的模板，缺乏专门针对\LaTeX 的开题报告模板，导致许多研究生在撰写过程中面临诸多挑战。

本课题旨在设计并实现一套适用于上海创智学院博士生的开题报告\LaTeX 模板。通过分析现有开题报告的结构和格式要求，结合\LaTeX 的排版优势，开发出一套功能完善、易于使用的模板。围绕这一研究方向，本课题从三个方面展开，分别探索了工作一、工作二和工作三的研究方法与技术路线，取得了初步的实验成果。在未来工作中，本课题拟从方面一、方面二、方面三继续深化研究，进一步完善模板功能，提高用户体验。

\vspace{3em}
\noindent{\bfseries 关键词：}\LaTeX，开题报告，模板设计，学术写作，排版自动化

\clearpage